\begin{afterword}
\summaryhead

The post continues ``Faster Than Science'': individual scientists can rationally believe things that Science has not yet confirmed. The easy case is a regularity spotted in existing data. The hard case is Einstein, who reached General Relativity with almost no data. From the principle that you cannot tell your speed inside a closed room, he got special relativity. From the aim of doing the same for acceleration, and from the equality of inertial and gravitational mass, he got curved spacetime. Mercury's orbit came out at the end as a check.

In an alternate history, the author says, the same theory would have needed decades of corrections to Newton before anyone could guess it. In Bayesian terms Einstein's feat was still induction: he learned the ``character of physical law'' from other laws and used it to predict a new one. His confidence before the eclipse, and against Bohr on quantum theory, was therefore justified. The post warns that history remembers the winners, advises ordinary reasoners to use every scrap of evidence, recommends evolutionary psychology as a model of reasoning from little data, and estimates that thousands of potential Einsteins are alive today.

\responsehead

The post's picture of Einstein's method is close to Einstein's own. In 1919 he described relativity as a theory built from general principles that any particular law must satisfy, and in 1952 he wrote that the field equations would be convincing even without the three classic tests. The post's warning that history remembers only the winners is also its own, and it answers a weakness of ``Faster Than Science'' the day before. The problems are in the details of the history, and each of them bears on the lesson.

Mercury was not only a final check. In 1913 Einstein computed it for his earlier Entwurf (``outline'') theory, which he held until 1915, and got too little, and in 1915 he named that failure as one of three reasons for dropping the theory. The claim that Einstein made no mistake requiring new data holds as worded. The chain did include a two-year wrong theory, held with the kind of certainty the post defends.

The principles did not all survive. The relativity of acceleration, which the post presents as a guiding property, is what Einstein named Mach's principle, and most experts polled in 1993 said general relativity does not fully satisfy it. The post also gives that name to the principle behind special relativity, which is the principle of relativity.

The eclipse is told in its skeptical version. The claim that the 1919 result favored Einstein ``essentially by luck'' is contested: a 1979 remeasurement of the plates supports the original conclusion that the Newtonian value was ruled out. And on Brush's reading of the reception, physicists gave the eclipse no more weight than Mercury, which does not fit the post's account of Science waiting for a novel prediction.

The Bayesian account adds a vocabulary, not a test. It repeats the argument of ``Einstein's Arrogance,'' which works only once the theory is known to be right, and the rule ``Wait for an occasion where they are wrong'' passes over the 1914 occasion. The claim that evolutionary psychologists later find out whether they were right is hedged and given no example.

In short: an account of Einstein's method close to his own, told through a simplified history that leaves out the Entwurf theory, and in which Mercury, Mach's principle and the eclipse fit the lesson better than the record does.
\end{afterword}
